\centerline{\bf \Huge Степень вхождения}

\begin{flushright}
        \scriptsize{Эпизод 7. Созданные человеком}
\end{flushright}


\textbf{Определение.} \textit{Степенью вхождения} простого числа $p$ в натуральное число $n$ будем называть наибольшее такое $k$, что $n$ делится на $p^k$. Обозначать для краткости будем $\nu_p(n)$ (это греческая буква "ню")


\problem{1}
Натуральные числа $a, b, c$ таковы, что число $\frac{a}{b}+\frac{b}{c}+\frac{c}{a}$ является целым. Верно ли, что $a b c-$ точный куб?


\problem{2}
Взаимно простые в совокупности натуральные числа $a, b, c$ удовлетворяют условию $a b=a c+b c$. Докажите, что $a b c$ - точный квадрат.

\problem{3}
\textbf{Формула Лежандра.} Докажите, что

$$
\nu_p(n!)=\left\lfloor\frac{n}{p}\right\rfloor+\left\lfloor\frac{n}{p^2}\right\rfloor+\left\lfloor\frac{n}{p^3}\right\rfloor+\ldots
$$

\problem{4} Докажите, что для любых натуральных $a_1, a_2, \ldots a_k$ число

$$
\frac{\left(a_1+a_2+\ldots+a_k\right)!}{a_{1}!a_{2}!\ldots a_{k}!}
$$


целое. Это число называется \textit{мультиномиальным коэффициентом}.

\problem{5} (a) Докажите, что наименьшее общее кратное чисел от $n$ до $2 n+1$ делится на $\frac{(2 n+1)!}{n!n!}$.

(b) Докажите, что наименьшее общее кратное чисел от 1 до $2 n+1$ больше $4^n$.

(c) Предположим, что среди чисел от 1 до $2 n+1$ ровно $k$ простых. Докажите, что $(2 n+1)^k>4^n$.

\problem{6}
Веня выписала в строчку 1000 натуральных чисел и обнаружил, что произведение любых двух соседних чисел является точным кубом. Докажите, что произведение двух крайних чисел тоже является точным кубом.

\problem{7}
Эля выписала на доску 57 различных натуральных чисел, не превосходящих 2024. Неля заметила, что среди любых трёх из этих чисел можно выбрать два, которые в произведении дадут точный квадрат. Докажите, что среди чисел, выписанных Элей на доску, встретится точный квадрат.